\chapter{Sliding Window Problems}
\chapterepigraph{A sliding window turns a quadratic ``look at every
subarray'' into a linear ``maintain an answer as the window moves.'' The
whole craft is choosing a data structure that supports the three window
operations cheaply: add the entering element, remove the leaving element,
and report the current answer. Pick the right structure and each of these
problems is the same loop with a different helper.}

\section{The Add / Remove / Report Template}
\label{sec:cses-sw-intro}

Every fixed-size sliding-window problem here processes windows of width
$k$ by advancing one step at a time: the right end gains an element, the
left end loses one, and we output the window's answer. The only thing that
varies is \emph{which structure} maintains that answer under those two
mutations.

\begin{ideabox}[Match the Answer to a Structure That Supports Removal]
\begin{itemize}
  \item \textbf{Sum / XOR} -- invertible: add the new, subtract/xor out
        the old in $O(1)$.
  \item \textbf{Min / Max} -- a \textbf{monotonic deque} gives $O(1)$
        amortised.
  \item \textbf{OR / bit aggregates} -- keep per-bit counts, since OR is
        not invertible.
  \item \textbf{Distinct count / mode / mex} -- a hash map of frequencies,
        sometimes with ``count of counts'' buckets or a set of absent
        values.
  \item \textbf{Median / cost} -- two balanced multisets split at the
        median, with running sums.
  \item \textbf{Inversions} -- a Fenwick tree over value ranks,
        updated incrementally as elements enter and leave.
\end{itemize}
\end{ideabox}

\begin{patternbox}[Why Removal Is the Hard Part]
Adding an element to a window is usually easy; \emph{removing} the element
that falls off the left is what dictates the data structure. Sums and XORs
are invertible so removal is trivial; min/max are not, which is why they
need a monotonic deque; order statistics need a balanced multiset or a
Fenwick tree. Reading a problem, ask first: ``how do I \emph{undo} the
leftmost element?''
\end{patternbox}

The sections cover, in order: window sum, minimum, XOR, OR, distinct
values, mode, mex, median, cost, inversions, and the advertisement
(maximum) variant. Each states the structure, the three window operations,
a dry run, complete C++, and complexity.

\newpage

\section{Sliding Window Sum}
\label{sec:cses-sw-sum}

\problemstatement{Given an array of $n$ integers and a window size $k$,
output the sum of every contiguous window of size $k$ (there are $n-k+1$
of them).}

\begin{patternbox}
Sum is \textbf{invertible}, so the window answer updates in $O(1)$: when
the window slides one step, add the new right element and subtract the old
left element. This is the simplest instance of the add/remove/report
template.
\end{patternbox}

\subsection*{A running total}

Compute the sum of the first window directly. Then each slide changes the
window by exactly two elements: one enters on the right, one leaves on the
left. Because addition has an inverse, the new sum is
$\textit{sum}+a[\text{new}]-a[\text{old}]$ -- no re-summation. (Equivalently,
a prefix-sum array answers each window as $\textit{pre}[i+k]-\textit{pre}
[i]$.)

\begin{ideabox}[Invertible Aggregates Slide for Free]
Any aggregate with an inverse -- sum, XOR, product over a field -- lets the
window update in constant time by undoing the departing element and
applying the arriving one. Reach for this before any heavier structure.
\end{ideabox}

\subsection*{Algorithm}

\begin{enumerate}
  \item Sum the first $k$ elements.
  \item For each subsequent window, add $a[i]$ and subtract $a[i-k]$;
        record the sum.
\end{enumerate}

\subsection*{Dry run}

\dryrun{$a=[1,3,2,5,1]$, $k=3$.}

\begin{itemize}
  \item Window $[1,3,2]$: sum $6$.
  \item Slide: $+5-1\Rightarrow[3,2,5]$ sum $10$.
  \item Slide: $+1-3\Rightarrow[2,5,1]$ sum $8$. Output $6,10,8$.
\end{itemize}

\subsection*{C++ implementation}

\begin{lstlisting}
#include <bits/stdc++.h>
using namespace std;

int main() {
    int n, k;
    cin >> n >> k;
    vector<long long> a(n);
    for (auto& x : a) cin >> x;

    long long sum = 0;
    for (int i = 0; i < k; i++) sum += a[i];
    cout << sum << " ";
    for (int i = k; i < n; i++) {
        sum += a[i] - a[i - k];               // add new, drop old
        cout << sum << " ";
    }
    cout << "\n";
    return 0;
}
\end{lstlisting}

\begin{complexitybox}
\textbf{Time Complexity.} $O(n)$ -- one pass with $O(1)$ per window.
\textbf{Space Complexity.} $O(1)$ beyond the input.
\end{complexitybox}

\begin{takeawaybox}
When the window statistic is invertible, sliding is a constant-time
add-and-subtract. Establish this baseline first; the harder window problems
are exactly those where the statistic (min, mode, median) has no cheap
inverse and demands a supporting structure.
\end{takeawaybox}

\section{Sliding Window Minimum}
\label{sec:cses-sw-minimum}

\problemstatement{Given an array of $n$ integers and a window size $k$,
output the minimum of every contiguous window of size $k$.}

\begin{patternbox}
Minimum has no cheap inverse, so a running value fails. The right tool is a
\textbf{monotonic deque} that keeps candidate minima in increasing order:
the front is always the current window's minimum, and each element is
pushed and popped at most once, giving $O(n)$ overall.
\end{patternbox}

\subsection*{The monotonic deque}

Maintain a deque of \emph{indices} whose values are strictly increasing
from front to back. When a new element arrives, pop from the back every
index whose value is $\ge$ the new value -- those can never be the minimum
again while the newer, smaller element is in the window. Then push the new
index. Pop from the front any index that has slid out of the window
(index $\le i-k$). The front index's value is the window minimum.

\begin{ideabox}[Discard Dominated Candidates]
An older element that is larger than a newer one is ``dominated'': as long
as the newer element is present, the older can never be the minimum, so it
is removed from the back forever. Each index enters and leaves the deque
once, which is why the total work is linear despite the inner popping.
\end{ideabox}

\subsection*{Algorithm}

\begin{enumerate}
  \item For each index $i$: pop back while its value $\ge a[i]$; push $i$.
  \item Pop front while it is $\le i-k$ (out of window).
  \item Once $i\ge k-1$, output $a[\text{front}]$.
\end{enumerate}

\subsection*{Dry run}

\dryrun{$a=[2,1,3,4,1]$, $k=2$.}

\begin{itemize}
  \item $i{=}0$: deque $[0]$. $i{=}1$: $a[1]{=}1\le a[0]$, pop $0$, push
        $1$; window min $a[1]=1$.
  \item $i{=}2$: push $2$ (deque $[1,2]$, but $1$ now out at $i-k=0$?
        front $1>0$ stays); min $a[1]=1$. Then front $1\le2-2=0$? no.
  \item Continuing yields minima $1,1,3,1$ for the four windows.
\end{itemize}

\subsection*{C++ implementation}

\begin{lstlisting}
#include <bits/stdc++.h>
using namespace std;

int main() {
    int n, k;
    cin >> n >> k;
    vector<int> a(n);
    for (int& x : a) cin >> x;

    deque<int> dq;                            // indices, values increasing
    for (int i = 0; i < n; i++) {
        while (!dq.empty() && a[dq.back()] >= a[i]) dq.pop_back();
        dq.push_back(i);
        if (dq.front() <= i - k) dq.pop_front();   // drop out-of-window
        if (i >= k - 1) cout << a[dq.front()] << " ";
    }
    cout << "\n";
    return 0;
}
\end{lstlisting}

\begin{complexitybox}
\textbf{Time Complexity.} $O(n)$ -- each index is pushed and popped at most
once. \textbf{Space Complexity.} $O(k)$ for the deque.
\end{complexitybox}

\begin{takeawaybox}
The monotonic deque is the canonical answer to sliding min/max: keep only
non-dominated candidates in monotonic order, front is the answer. It is the
single most reused sliding-window structure and appears again as the
maximum (advertisement) variant.
\end{takeawaybox}

\section{Sliding Window XOR}
\label{sec:cses-sw-xor}

\problemstatement{Given an array of $n$ integers and a window size $k$,
output the bitwise XOR of every contiguous window of size $k$.}

\begin{patternbox}
XOR is its own inverse ($x\oplus x=0$), so, like sum, the window XOR slides
in $O(1)$: XOR in the entering element and XOR out the leaving one.
Equivalently, a \textbf{prefix-XOR} array answers each window as
$\textit{pre}[i+k]\oplus\textit{pre}[i]$.
\end{patternbox}

\subsection*{Prefix XOR}

Let $\textit{pre}[0]=0$ and $\textit{pre}[i]=a[0]\oplus\cdots\oplus
a[i-1]$. Then the XOR of $a[l..r]$ is $\textit{pre}[r+1]\oplus
\textit{pre}[l]$, because the shared prefix $a[0..l-1]$ cancels itself
(each bit XORed twice returns to $0$). A window of size $k$ starting at $i$
is $\textit{pre}[i+k]\oplus\textit{pre}[i]$.

\begin{ideabox}[Self-Inverse Means Cancellation]
Because $x\oplus x=0$, overlapping prefixes cancel exactly, so a
prefix-XOR array reduces any window XOR to a single operation. The same
self-inverse property lets a running XOR drop the departing element by
simply XORing it again.
\end{ideabox}

\subsection*{Algorithm}

\begin{enumerate}
  \item Build the prefix-XOR array (or maintain a running XOR).
  \item For each window $i$, output
        $\textit{pre}[i+k]\oplus\textit{pre}[i]$.
\end{enumerate}

\subsection*{Dry run}

\dryrun{$a=[1,2,3,4]$, $k=2$.}

\begin{itemize}
  \item $\textit{pre}=[0,1,3,0,4]$.
  \item Window $[1,2]$: $\textit{pre}[2]\oplus\textit{pre}[0]=3$.
        $[2,3]$: $\textit{pre}[3]\oplus\textit{pre}[1]=0\oplus1=1$.
        $[3,4]$: $\textit{pre}[4]\oplus\textit{pre}[2]=4\oplus3=7$.
  \item Output $3,1,7$.
\end{itemize}

\subsection*{C++ implementation}

\begin{lstlisting}
#include <bits/stdc++.h>
using namespace std;

int main() {
    int n, k;
    cin >> n >> k;
    vector<int> a(n);
    for (int& x : a) cin >> x;

    vector<int> pre(n + 1, 0);
    for (int i = 0; i < n; i++) pre[i + 1] = pre[i] ^ a[i];

    for (int i = 0; i + k <= n; i++)
        cout << (pre[i + k] ^ pre[i]) << " ";     // window XOR
    cout << "\n";
    return 0;
}
\end{lstlisting}

\begin{complexitybox}
\textbf{Time Complexity.} $O(n)$. \textbf{Space Complexity.} $O(n)$ for the
prefix array (or $O(1)$ with a running XOR).
\end{complexitybox}

\begin{takeawaybox}
XOR joins sum in the ``invertible, slides in $O(1)$'' category thanks to
being its own inverse. Prefix XOR is the range-query form and reappears in
counting subarrays with a target XOR. Contrast with OR (next), which is
\emph{not} invertible and needs per-bit counts.
\end{takeawaybox}

\section{Sliding Window OR}
\label{sec:cses-sw-or}

\problemstatement{Given an array of $n$ integers and a window size $k$,
output the bitwise OR of every contiguous window of size $k$.}

\begin{patternbox}
OR is \textbf{not invertible}: removing an element cannot simply undo its
bits, because another element in the window may still set them. The fix is
to track, for each bit position, \emph{how many} window elements set that
bit -- then the OR has a bit set exactly when its count is positive.
\end{patternbox}

\subsection*{Per-bit counts}

Maintain an array $\textit{cnt}[b]$ = number of elements in the current
window whose bit $b$ is set. Adding an element increments $\textit{cnt}[b]$
for each of its set bits; removing an element decrements them. The OR of
the window has bit $b$ set iff $\textit{cnt}[b]>0$. Reconstruct the OR by
summing $2^b$ over all bits with positive count. This restores $O(1)$
amortised updates (constant number of bits) despite OR's lack of an
inverse.

\begin{ideabox}[Count Contributors per Bit]
When an aggregate is not invertible, replace ``does any element set this
bit?'' with ``how many do?'' A count is invertible (increment/decrement),
and the OR is recovered by testing each count against zero. The same idea
handles AND (count zeros) and other non-invertible bitwise folds.
\end{ideabox}

\subsection*{Algorithm}

\begin{enumerate}
  \item For the first window, increment $\textit{cnt}[b]$ for every set
        bit of each element.
  \item To slide: add the new element's bits, remove the old element's
        bits.
  \item Output $\sum_{b:\,\textit{cnt}[b]>0}2^b$ for each window.
\end{enumerate}

\subsection*{Dry run}

\dryrun{$a=[1,2,3]$ (bits: $1{=}001$, $2{=}010$, $3{=}011$), $k=2$.}

\begin{itemize}
  \item Window $[1,2]$: bit$0$ count $1$, bit$1$ count $1$ $\Rightarrow$
        OR $=3$.
  \item Slide to $[2,3]$: remove $1$ (bit$0\to0$), add $3$ (bit$0\to1$,
        bit$1\to2$) $\Rightarrow$ bits $0,1$ set, OR $=3$.
  \item Output $3,3$.
\end{itemize}

\subsection*{C++ implementation}

\begin{lstlisting}
#include <bits/stdc++.h>
using namespace std;

int main() {
    int n, k;
    cin >> n >> k;
    vector<int> a(n);
    for (int& x : a) cin >> x;

    const int BITS = 31;
    vector<int> cnt(BITS, 0);
    auto add = [&](int x, int delta) {
        for (int b = 0; b < BITS; b++) if (x & (1 << b)) cnt[b] += delta;
    };
    auto currentOR = [&]() {
        int res = 0;
        for (int b = 0; b < BITS; b++) if (cnt[b] > 0) res |= (1 << b);
        return res;
    };

    for (int i = 0; i < k; i++) add(a[i], +1);
    cout << currentOR() << " ";
    for (int i = k; i < n; i++) {
        add(a[i], +1); add(a[i - k], -1);         // add new, remove old
        cout << currentOR() << " ";
    }
    cout << "\n";
    return 0;
}
\end{lstlisting}

\begin{complexitybox}
\textbf{Time Complexity.} $O(n\cdot\text{BITS})$ -- constant $31$ bits per
element. \textbf{Space Complexity.} $O(\text{BITS})$.
\end{complexitybox}

\begin{takeawaybox}
Non-invertible aggregates like OR are made sliding-friendly by counting
contributors per bit, so removal becomes a decrement. This
``count-then-threshold'' conversion is the general recipe whenever an
operation lacks an inverse.
\end{takeawaybox}

\section{Sliding Window Distinct Values}
\label{sec:cses-sw-distinct}

\problemstatement{Given an array of $n$ integers and a window size $k$,
output the number of distinct values in every contiguous window of size
$k$.}

\begin{patternbox}
Distinct-count slides cleanly with a \textbf{frequency hash map}: track how
many times each value appears in the window. A value contributes to the
distinct count exactly while its frequency is positive, so the count
changes only when a frequency crosses between $0$ and $1$.
\end{patternbox}

\subsection*{Frequency map with a distinct counter}

Keep $\textit{freq}[v]$ and a running $\textit{distinct}$. Adding a value
$v$: if its frequency was $0$, it is newly present, so increment
$\textit{distinct}$; then increment $\textit{freq}[v]$. Removing a value
$v$: decrement $\textit{freq}[v]$; if it drops to $0$, the value has left
the window, so decrement $\textit{distinct}$. Each slide performs one add
and one remove, both $O(1)$ amortised.

\begin{ideabox}[Watch the Zero Boundary]
The distinct count only moves when a frequency transitions $0\to1$ (a value
appears) or $1\to0$ (a value disappears). Checking that boundary on each
add/remove keeps the counter exact without rescanning the window.
\end{ideabox}

\subsection*{Algorithm}

\begin{enumerate}
  \item Add the first $k$ elements, updating $\textit{freq}$ and
        $\textit{distinct}$.
  \item For each slide: add $a[i]$ (increment distinct if it was absent),
        remove $a[i-k]$ (decrement distinct if it becomes absent).
  \item Output $\textit{distinct}$ per window.
\end{enumerate}

\subsection*{Dry run}

\dryrun{$a=[1,2,1,3]$, $k=3$.}

\begin{itemize}
  \item Window $[1,2,1]$: freqs $\{1{:}2,2{:}1\}$, distinct $=2$.
  \item Slide to $[2,1,3]$: remove a $1$ (freq $2{\to}1$, still present),
        add $3$ (new) $\Rightarrow$ distinct $=3$.
  \item Output $2,3$.
\end{itemize}

\subsection*{C++ implementation}

\begin{lstlisting}
#include <bits/stdc++.h>
using namespace std;

int main() {
    int n, k;
    cin >> n >> k;
    vector<int> a(n);
    for (int& x : a) cin >> x;

    unordered_map<int,int> freq;
    int distinct = 0;
    auto add = [&](int v) { if (freq[v]++ == 0) distinct++; };
    auto remove = [&](int v) { if (--freq[v] == 0) distinct--; };

    for (int i = 0; i < k; i++) add(a[i]);
    cout << distinct << " ";
    for (int i = k; i < n; i++) {
        add(a[i]); remove(a[i - k]);
        cout << distinct << " ";
    }
    cout << "\n";
    return 0;
}
\end{lstlisting}

\begin{complexitybox}
\textbf{Time Complexity.} $O(n)$ amortised with a hash map.
\textbf{Space Complexity.} $O(k)$ distinct keys at most.
\end{complexitybox}

\begin{takeawaybox}
Distinct-value counting is a frequency map plus a counter that only reacts
at the zero boundary. This add/remove-with-boundary-check idiom generalises
to ``count of values with property P'' whenever P depends on a per-value
frequency.
\end{takeawaybox}

\section{Sliding Window Mode}
\label{sec:cses-sw-mode}

\problemstatement{Given an array of $n$ integers and a window size $k$,
output the mode of every contiguous window of size $k$ -- the value
occurring most often (if several tie, the smallest such value).}

\begin{patternbox}
The mode must survive both additions and \emph{removals}, so a single
``max frequency'' is not enough. Keep frequency buckets: for each frequency
$f$, an ordered set of the values that currently occur $f$ times. The mode
is the smallest value in the highest non-empty bucket.
\end{patternbox}

\subsection*{Frequency buckets of values}

Maintain $\textit{freq}[v]$ and $\textit{bucket}[f]$ = an ordered set of
values whose current frequency is exactly $f$, plus $\textit{maxFreq}$.
Adding a value $v$ moves it from $\textit{bucket}[\textit{freq}[v]]$ to
$\textit{bucket}[\textit{freq}[v]+1]$ and may raise $\textit{maxFreq}$.
Removing $v$ moves it down a bucket; if that empties the top bucket,
$\textit{maxFreq}$ decreases. The mode is the smallest value in
$\textit{bucket}[\textit{maxFreq}]$, read in $O(\log)$ from the ordered
set's begin.

\begin{ideabox}[Buckets Make the Mode Removal-Safe]
The trouble with modes under removal is that the top frequency can drop.
Grouping values by their exact frequency in ordered buckets lets a removal
adjust $\textit{maxFreq}$ in $O(1)$ and still expose the smallest tied
value instantly -- something a lone counter cannot do.
\end{ideabox}

\subsection*{Algorithm}

\begin{enumerate}
  \item Maintain $\textit{freq}$, $\textit{bucket}[f]$ (ordered set of
        values), and $\textit{maxFreq}$.
  \item On add/remove of $v$: erase $v$ from its old bucket, change its
        frequency, insert into the new bucket, and update
        $\textit{maxFreq}$.
  \item The mode is $\ast\textit{bucket}[\textit{maxFreq}].\text{begin}()$;
        output it per window.
\end{enumerate}

\subsection*{Dry run}

\dryrun{$a=[1,2,2,1,1]$, $k=3$.}

\begin{itemize}
  \item $[1,2,2]$: freqs $1{:}1,2{:}2$; maxFreq $2$, bucket$[2]=\{2\}$;
        mode $2$.
  \item Slide to $[2,2,1]$: still $2{:}2,1{:}1$; mode $2$.
  \item Slide to $[2,1,1]$: $1{:}2,2{:}1$; bucket$[2]=\{1\}$; mode $1$.
        Output $2,2,1$.
\end{itemize}

\subsection*{C++ implementation}

\begin{lstlisting}
#include <bits/stdc++.h>
using namespace std;

int main() {
    int n, k;
    cin >> n >> k;
    vector<int> a(n);
    for (int& x : a) cin >> x;

    unordered_map<int,int> freq;
    map<int, set<int>> bucket;                // frequency -> values
    int maxFreq = 0;

    auto change = [&](int v, int delta) {
        int f = freq[v];
        if (f > 0) bucket[f].erase(v);
        if (f > 0 && bucket[f].empty()) bucket.erase(f);
        f += delta;
        freq[v] = f;
        if (f > 0) bucket[f].insert(v);
        maxFreq = bucket.empty() ? 0 : bucket.rbegin()->first;
    };

    for (int i = 0; i < k; i++) change(a[i], +1);
    cout << *bucket[maxFreq].begin() << " ";
    for (int i = k; i < n; i++) {
        change(a[i], +1); change(a[i - k], -1);
        cout << *bucket[maxFreq].begin() << " ";
    }
    cout << "\n";
    return 0;
}
\end{lstlisting}

\begin{complexitybox}
\textbf{Time Complexity.} $O(n\log k)$ -- each add/remove touches ordered
sets. \textbf{Space Complexity.} $O(k)$.
\end{complexitybox}

\begin{takeawaybox}
Tracking the mode under removals needs frequency-to-values buckets so the
top frequency and its smallest value are both maintainable. The ``count of
counts'' idea -- indexing by frequency, not by value -- is the general trick
for order statistics on frequencies.
\end{takeawaybox}

\section{Sliding Window Mex}
\label{sec:cses-sw-mex}

\problemstatement{Given an array of $n$ non-negative integers and a window
size $k$, output the mex of every contiguous window -- the smallest
non-negative integer \emph{not} present in the window.}

\begin{patternbox}
The mex of a size-$k$ window is at most $k$ (with $k$ elements, some value
in $0\ldots k$ must be missing). So track only values in $[0,k]$ with a
frequency map and keep an ordered \textbf{set of absent values}; the mex is
its smallest element.
\end{patternbox}

\subsection*{A set of the missing small values}

Values larger than $k$ can never be the mex, so ignore them. Initialise a
set $\textit{absent}=\{0,1,\ldots,k\}$ and a frequency map. When a value
$v\le k$ enters and its frequency goes $0\to1$, erase $v$ from
$\textit{absent}$ (now present). When it leaves and its frequency returns
to $0$, insert $v$ back into $\textit{absent}$. The mex is the smallest
element of $\textit{absent}$ -- read in $O(\log k)$ from its begin.

\begin{ideabox}[Bound the Universe, Then Track Absence]
Because the mex cannot exceed the window size, the relevant universe is
just $[0,k]$. Maintaining which of those are absent -- rather than which are
present -- puts the answer (the smallest absent value) at the front of an
ordered set.
\end{ideabox}

\subsection*{Algorithm}

\begin{enumerate}
  \item $\textit{absent}=\{0,\ldots,k\}$; empty frequency map.
  \item Add the first $k$ elements (erasing from $\textit{absent}$ on
        $0\to1$).
  \item Each slide: add $a[i]$, remove $a[i-k]$ (inserting into
        $\textit{absent}$ on $1\to0$); output $\ast\textit{absent}.
        \text{begin}()$.
\end{enumerate}

\subsection*{Dry run}

\dryrun{$a=[0,1,3,0]$, $k=3$.}

\begin{itemize}
  \item $\textit{absent}=\{0,1,2,3\}$. Window $[0,1,3]$: erase $0,1,3$;
        absent $=\{2\}$; mex $2$.
  \item Slide to $[1,3,0]$: remove first $0$ (freq $0$, reinsert $0$),
        add $0$ (erase $0$ again); absent $=\{2\}$; mex $2$.
  \item Output $2,2$.
\end{itemize}

\subsection*{C++ implementation}

\begin{lstlisting}
#include <bits/stdc++.h>
using namespace std;

int main() {
    int n, k;
    cin >> n >> k;
    vector<int> a(n);
    for (int& x : a) cin >> x;

    set<int> absent;
    for (int v = 0; v <= k; v++) absent.insert(v);
    unordered_map<int,int> freq;

    auto add = [&](int v) {
        if (v <= k && freq[v]++ == 0) absent.erase(v);
    };
    auto remove = [&](int v) {
        if (v <= k && --freq[v] == 0) absent.insert(v);
    };

    for (int i = 0; i < k; i++) add(a[i]);
    cout << *absent.begin() << " ";
    for (int i = k; i < n; i++) {
        add(a[i]); remove(a[i - k]);
        cout << *absent.begin() << " ";
    }
    cout << "\n";
    return 0;
}
\end{lstlisting}

\begin{complexitybox}
\textbf{Time Complexity.} $O(n\log k)$ -- ordered-set operations bounded by
the $[0,k]$ universe. \textbf{Space Complexity.} $O(k)$.
\end{complexitybox}

\begin{takeawaybox}
Mex problems shrink to a bounded universe (here $[0,k]$) and then track
\emph{absence} in an ordered set so the answer is the set's minimum. Bounding
the universe before choosing a structure is the recurring first move for
mex.
\end{takeawaybox}

\section{Sliding Window Median}
\label{sec:cses-sw-median}

\problemstatement{Given an array of $n$ integers and a window size $k$,
output the median of every contiguous window (for even $k$, the lower of
the two middle values).}

\begin{patternbox}
The median needs the middle element under both insertion and deletion. Two
balanced \textbf{multisets} -- a \emph{low} half and a \emph{high} half --
split at the median do exactly this: keep the low half's size equal to (or
one more than) the high half, and the median is the low half's maximum.
\end{patternbox}

\subsection*{Two balanced halves}

Maintain \textit{low} (a multiset of the smaller elements) and \textit{high}
(the larger elements) with the invariant $|\textit{low}|=\lceil m/2\rceil$
for a window of size $m$. Insertion routes the new value to the correct half
and then rebalances by moving one element across if sizes drift. Deletion
(when the left element leaves) erases it from whichever half holds it -- a
multiset supports erase-by-value in $O(\log k)$ -- and rebalances. The
median is always $\ast\textit{low}.\text{rbegin}()$ (its largest element).

\begin{ideabox}[Balanced Multisets Give Erasable Heaps]
Ordinary heaps cannot delete an arbitrary element, which sliding windows
require. Two multisets behave like two heaps that \emph{also} support
$O(\log k)$ removal of a specific value, keeping the median at the boundary
between the halves.
\end{ideabox}

\subsection*{Algorithm}

\begin{enumerate}
  \item Insert into \textit{low} or \textit{high} by comparing with the
        current median; rebalance so $|\textit{low}|-|\textit{high}|
        \in\{0,1\}$.
  \item On slide, erase the departing value from its half and rebalance.
  \item Output $\ast\textit{low}.\text{rbegin}()$ per window.
\end{enumerate}

\subsection*{Dry run}

\dryrun{$a=[2,4,1,3]$, $k=3$.}

\begin{itemize}
  \item $[2,4,1]$ sorted $1,2,4$: low $=\{1,2\}$, high $=\{4\}$; median
        $2$.
  \item Slide to $[4,1,3]$ sorted $1,3,4$: erase $2$, insert $3$; low
        $=\{1,3\}$, high $=\{4\}$; median $3$.
  \item Output $2,3$.
\end{itemize}

\subsection*{C++ implementation}

\begin{lstlisting}
#include <bits/stdc++.h>
using namespace std;

int main() {
    int n, k;
    cin >> n >> k;
    vector<int> a(n);
    for (int& x : a) cin >> x;

    multiset<int> low, high;                  // low: smaller half
    auto rebalance = [&]() {
        while (low.size() > high.size() + 1) {
            high.insert(*low.rbegin());
            low.erase(prev(low.end()));
        }
        while (low.size() < high.size()) {
            low.insert(*high.begin());
            high.erase(high.begin());
        }
    };
    auto insert = [&](int x) {
        if (low.empty() || x <= *low.rbegin()) low.insert(x);
        else high.insert(x);
        rebalance();
    };
    auto erase = [&](int x) {
        auto it = low.find(x);
        if (it != low.end()) low.erase(it);
        else high.erase(high.find(x));
        rebalance();
    };

    for (int i = 0; i < k; i++) insert(a[i]);
    cout << *low.rbegin() << " ";
    for (int i = k; i < n; i++) {
        insert(a[i]); erase(a[i - k]);
        cout << *low.rbegin() << " ";
    }
    cout << "\n";
    return 0;
}
\end{lstlisting}

\begin{complexitybox}
\textbf{Time Complexity.} $O(n\log k)$ -- each insert/erase/rebalance is
$O(\log k)$. \textbf{Space Complexity.} $O(k)$.
\end{complexitybox}

\begin{takeawaybox}
The two-balanced-multisets structure keeps the median at the split point
under both insertion and deletion -- the erasable-heap the sliding window
needs. It directly powers the next problem, sliding window cost, by adding
running sums to each half.
\end{takeawaybox}

\section{Sliding Window Cost}
\label{sec:cses-sw-cost}

\problemstatement{Given an array of $n$ integers and a window size $k$, for
every window output the minimum total cost of making all its elements
equal, where the cost of changing a value by $d$ is $\lvert d\rvert$. The
optimal target is the window's median.}

\begin{patternbox}
Minimising $\sum\lvert a_i-x\rvert$ is achieved at $x=\text{median}$, so
this is the median structure of the previous section \emph{plus running
sums}. With the two halves' sizes and sums maintained, the cost is a
closed-form expression evaluated in $O(1)$ per window.
\end{patternbox}

\subsection*{Cost from the two halves}

Reuse the low/high multisets that keep the median at the boundary, and also
maintain $\textit{lowSum}$ and $\textit{highSum}$. With median $x$, the
elements in \textit{low} are all $\le x$ and cost
$x\cdot|\textit{low}|-\textit{lowSum}$ to raise; those in \textit{high} are
$\ge x$ and cost $\textit{highSum}-x\cdot|\textit{high}|$ to lower. The
total is
\[
  \text{cost}=\big(x\,|\textit{low}|-\textit{lowSum}\big)+
  \big(\textit{highSum}-x\,|\textit{high}|\big).
\]
Every insert/erase updates the corresponding sum alongside the multiset.

\begin{ideabox}[The Median Minimises Absolute Deviation]
For $\sum\lvert a_i-x\rvert$, moving $x$ toward the median always reduces
the total until $x$ reaches it -- a standard exchange argument. So the only
statistic needed is the median plus the sums of the two halves, all of
which the balanced multisets already track.
\end{ideabox}

\subsection*{Algorithm}

\begin{enumerate}
  \item Maintain low/high multisets (median at the split) and their sums.
  \item On insert/erase, update sizes and sums.
  \item Per window, compute the closed-form cost with $x=\ast\textit{low}.
        \text{rbegin}()$.
\end{enumerate}

\subsection*{Dry run}

\dryrun{Window $[1,4,3]$, median $3$.}

\begin{itemize}
  \item low $=\{1,3\}$ (sum $4$), high $=\{4\}$ (sum $4$), $x=3$.
  \item cost $=(3\cdot2-4)+(4-3\cdot1)=2+1=3$ -- raise $1\to3$ (cost $2$),
        lower $4\to3$ (cost $1$).
\end{itemize}

\subsection*{C++ implementation}

\begin{lstlisting}
#include <bits/stdc++.h>
using namespace std;

int main() {
    int n, k;
    cin >> n >> k;
    vector<long long> a(n);
    for (auto& x : a) cin >> x;

    multiset<long long> low, high;
    long long lowSum = 0, highSum = 0;
    auto rebalance = [&]() {
        while (low.size() > high.size() + 1) {
            long long v = *low.rbegin(); low.erase(prev(low.end()));
            lowSum -= v; high.insert(v); highSum += v;
        }
        while (low.size() < high.size()) {
            long long v = *high.begin(); high.erase(high.begin());
            highSum -= v; low.insert(v); lowSum += v;
        }
    };
    auto insert = [&](long long x) {
        if (low.empty() || x <= *low.rbegin()) { low.insert(x); lowSum += x; }
        else { high.insert(x); highSum += x; }
        rebalance();
    };
    auto erase = [&](long long x) {
        auto it = low.find(x);
        if (it != low.end()) { low.erase(it); lowSum -= x; }
        else { high.erase(high.find(x)); highSum -= x; }
        rebalance();
    };
    auto cost = [&]() {
        long long x = *low.rbegin();
        return (x * (long long)low.size() - lowSum) +
               (highSum - x * (long long)high.size());
    };

    for (int i = 0; i < k; i++) insert(a[i]);
    cout << cost() << " ";
    for (int i = k; i < n; i++) {
        insert(a[i]); erase(a[i - k]);
        cout << cost() << " ";
    }
    cout << "\n";
    return 0;
}
\end{lstlisting}

\begin{complexitybox}
\textbf{Time Complexity.} $O(n\log k)$ -- median maintenance dominates.
\textbf{Space Complexity.} $O(k)$.
\end{complexitybox}

\begin{takeawaybox}
Minimising absolute deviation is a median problem; carrying the two halves'
sums alongside the median turns the cost into a one-line formula. Augmenting
an order-statistic structure with aggregate sums is a broadly useful move.
\end{takeawaybox}

\section{Sliding Window Inversions}
\label{sec:cses-sw-inversions}

\problemstatement{Given an array of $n$ integers and a window size $k$,
output the number of inversions in every contiguous window -- pairs of
positions $i<j$ within the window with $a_i>a_j$.}

\begin{patternbox}
Recomputing inversions per window is $O(nk\log k)$. Instead maintain a
running inversion count and a \textbf{Fenwick tree} over value ranks:
adding the new right element and removing the old left element each change
the count by a range query answerable in $O(\log n)$.
\end{patternbox}

\subsection*{Incremental inversion count}

Keep a Fenwick tree of the value counts currently in the window and a
running total $\textit{inv}$. When a new element $x$ enters on the
\emph{right}, it forms an inversion with every current element strictly
greater than $x$ (all of which sit to its left), so
$\textit{inv}\mathrel{+}=(\text{window size})-\#\{\le x\}$; then insert
$x$. When the leftmost element $y$ leaves, it was to the left of everything
else, so it formed an inversion with every current element strictly less
than $y$: $\textit{inv}\mathrel{-}=\#\{<y\}$; then remove $y$. Coordinate-
compress values so ranks fit the Fenwick tree.

\begin{ideabox}[Add on the Right, Remove on the Left, Query the Rank]
The window's geometry makes the updates clean: a right-entering element
pairs with larger existing elements; a left-leaving element pairs with
smaller ones. Both counts are Fenwick prefix queries, so each slide is
$O(\log n)$.
\end{ideabox}

\subsection*{Algorithm}

\begin{enumerate}
  \item Compress values to ranks $1\ldots m$. Build the first window,
        adding elements and accumulating inversions.
  \item Per slide: remove the leftmost ($\textit{inv}\mathrel{-}=$ count
        of smaller present), then add the new ($\textit{inv}
        \mathrel{+}=$ count of larger present).
  \item Output $\textit{inv}$ per window.
\end{enumerate}

\subsection*{Dry run}

\dryrun{$a=[2,1,3]$, $k=2$.}

\begin{itemize}
  \item Window $[2,1]$: $2>1$ is one inversion $\Rightarrow1$.
  \item Slide to $[1,3]$: remove $2$ (it had $1$ smaller element $\to
        \textit{inv}=0$), add $3$ (nothing greater present $\to$ stays
        $0$).
  \item Output $1,0$.
\end{itemize}

\subsection*{C++ implementation}

\begin{lstlisting}
#include <bits/stdc++.h>
using namespace std;

struct BIT {
    int n; vector<int> t;
    BIT(int n) : n(n), t(n + 1, 0) {}
    void update(int i, int v) { for (; i <= n; i += i & -i) t[i] += v; }
    int query(int i) { int s = 0; for (; i > 0; i -= i & -i) s += t[i]; return s; }
};

int main() {
    int n, k;
    cin >> n >> k;
    vector<int> a(n);
    for (int& x : a) cin >> x;

    vector<int> vals(a);
    sort(vals.begin(), vals.end());
    vals.erase(unique(vals.begin(), vals.end()), vals.end());
    auto rankOf = [&](int x) {
        return int(lower_bound(vals.begin(), vals.end(), x) - vals.begin()) + 1;
    };
    int m = vals.size();
    BIT bit(m);

    long long inv = 0;
    int inWin = 0;
    for (int i = 0; i < k; i++) {              // build first window
        int r = rankOf(a[i]);
        inv += inWin - bit.query(r);           // existing elements > a[i]
        bit.update(r, +1); inWin++;
    }
    cout << inv << " ";
    for (int i = k; i < n; i++) {
        int ry = rankOf(a[i - k]);             // remove leftmost
        inv -= bit.query(ry - 1);              // present elements < y
        bit.update(ry, -1); inWin--;

        int rx = rankOf(a[i]);                 // add new right
        inv += inWin - bit.query(rx);          // present elements > x
        bit.update(rx, +1); inWin++;
        cout << inv << " ";
    }
    cout << "\n";
    return 0;
}
\end{lstlisting}

\begin{complexitybox}
\textbf{Time Complexity.} $O(n\log n)$ -- each slide does a constant number
of Fenwick operations. \textbf{Space Complexity.} $O(n)$ for the Fenwick
tree and compression.
\end{complexitybox}

\begin{takeawaybox}
Sliding-window inversions become linear-log by maintaining a running count
plus a Fenwick tree of present values: the right-entering element pairs
with larger elements, the left-leaving with smaller. Recognising which side
an element joins/leaves fixes which prefix query to run.
\end{takeawaybox}

\section{Sliding Window Advertisement (Maximum)}
\label{sec:cses-sw-advertisement}

\problemstatement{Given an array of $n$ integers and a window size $k$,
output the maximum of every contiguous window. (This ``advertisement''
variant is the maximum counterpart of the sliding window minimum.)}

\begin{patternbox}
Maximum, like minimum, is not invertible, so the same \textbf{monotonic
deque} applies -- kept in \emph{decreasing} order this time, so the front is
always the current window's maximum. Each index is pushed and popped once,
giving $O(n)$.
\end{patternbox}

\subsection*{A decreasing deque}

Maintain a deque of indices whose values are strictly decreasing from front
to back. When a new element arrives, pop from the back every index whose
value is $\le$ the new value -- those are dominated and can never be the
maximum again while the newer, larger element is present. Push the new
index, drop the front if it has left the window, and read the front as the
window maximum.

\begin{ideabox}[Mirror of the Minimum Deque]
The only change from sliding window minimum is the comparison direction:
pop while the back is $\le$ the incoming value (instead of $\ge$), keeping
the deque decreasing so its front is the maximum. One structure, two mirror
uses.
\end{ideabox}

\subsection*{Algorithm}

\begin{enumerate}
  \item For each index $i$: pop back while its value $\le a[i]$; push $i$.
  \item Pop front while it is $\le i-k$ (out of window).
  \item Once $i\ge k-1$, output $a[\text{front}]$.
\end{enumerate}

\subsection*{Dry run}

\dryrun{$a=[3,1,4,2]$, $k=2$.}

\begin{itemize}
  \item $[3,1]$: deque front $3$; max $3$.
  \item $[1,4]$: $4$ pops $1$ (and $3$ leaves window); max $4$.
  \item $[4,2]$: front $4$; max $4$. Output $3,4,4$.
\end{itemize}

\subsection*{C++ implementation}

\begin{lstlisting}
#include <bits/stdc++.h>
using namespace std;

int main() {
    int n, k;
    cin >> n >> k;
    vector<int> a(n);
    for (int& x : a) cin >> x;

    deque<int> dq;                            // indices, values decreasing
    for (int i = 0; i < n; i++) {
        while (!dq.empty() && a[dq.back()] <= a[i]) dq.pop_back();
        dq.push_back(i);
        if (dq.front() <= i - k) dq.pop_front();
        if (i >= k - 1) cout << a[dq.front()] << " ";
    }
    cout << "\n";
    return 0;
}
\end{lstlisting}

\begin{complexitybox}
\textbf{Time Complexity.} $O(n)$ -- each index enters and leaves the deque
once. \textbf{Space Complexity.} $O(k)$.
\end{complexitybox}

\begin{takeawaybox}
Sliding window maximum is the minimum deque mirrored: pop while the back is
no larger than the incoming value. Master the monotonic deque once and both
extremes -- and many ``nearest greater/smaller'' problems -- follow from it.
\end{takeawaybox}

